\documentclass{article}
\usepackage[utf8]{inputenc}
\usepackage[T1]{fontenc}
\usepackage{graphicx}
\usepackage{algorithm}
\usepackage{algpseudocode}

\usepackage{listings}
\usepackage{xcolor}
\usepackage[margin=2cm]{geometry}
\usepackage{float}

\usepackage{amssymb} %Ajoute des symboles mathématiques comme R, Z, forall, exists, etc....
\usepackage{amsmath}
\usepackage{amsthm}


\lstdefinestyle{monstyle}{
    backgroundcolor=\color{white},
    commentstyle=\color{gray},
    keywordstyle=\color{blue},
    numberstyle=\tiny\color{gray},
    basicstyle=\ttfamily\small,
    breaklines=true,
    numbers=left,
    numbersep=5pt,
    frame=single,
}
\lstset{style=monstyle}



\title{\textbf{TD n°5 : Initiation à l'algorithmique}}
\author{Derrez Lorenzo}

\begin{document}
\maketitle

\section*{\textbf{Exercice 0}}

% Insérer image 

\section*{\textbf{Exercice 1}}

\subsection*{Question 1:}
On obtient les tableaux suivant : 
\begin{itemize}
    \item $ T = [10, 20, 60, 50, 30, 40] $
    \item $ T = [10, 20, 60, 50, 30, 40] $
    \item $ T = [10, 20, 30, 50, 60, 40] $
    \item $ T = [10, 20, 30, 40, 60, 50] $
    \item $ T = [10, 20, 30, 40, 50, 60] $
    \item $ T = [10, 20, 30, 40, 50, 60] $
\end{itemize}

\subsection*{Question 2}

Jeu de 6 exemples pour \textit{indiceDuMinimum} :
\\
$
\begin{cases}
    & E : T = [1], i = 0
    \\
    & S : j = 0   
\end{cases}
$
\\
$
\begin{cases}
    & E : T = [1,1], i = 0
    \\
    & S : j = 0   
\end{cases}
$
\\
$
\begin{cases}
    & E : T = [3, 6, 2], i = 0
    \\
    & S : j = 2
\end{cases}
$
\\
$
\begin{cases}
    & E : T = [1, 6? 7, 2], i = 1
    \\
    & S : j = 3
\end{cases}
$
\\
$
\begin{cases}
    & E : T = [1, 6, 2, 2], i = 1
    \\
    & S : j = 2   
\end{cases}
$
\\
$
\begin{cases}
    & E : T = [12, 3, 10], i = 2
    \\
    & S : j = 2
\end{cases}
$

\subsection*{Question 3}

On définiti le problème \textit{IndiceDuMinimum} :

\begin{algorithm}[H]
    \caption{\textit{IndiceDuMinimum}}
    \begin{algorithmic}[1]
        \State E : un tableau $T$ d'entiers, et un entier i compris entre 0 et $ |T|-1$
        \State S : un entier j compris entre i et $ |T|-1$ tel que T[j] est le minimum entre i et |T|-1.
        \State def IndiceDuMinimum(i:int, T:Array[integer]) -> int:
        \State \qquad int n $\leftarrow$ size(T)
        \State \qquad int min\_ act $\leftarrow$ T[i]
        \State \qquad int ind\_ act $\leftarrow$ i 
        \State \qquad for k in range(i, n):
            \State \qquad \qquad if T[k] < min\_ act:
                \State \qquad \qquad \qquad min\_ act $\leftarrow$ T[k]
                \State \qquad \qquad \qquad ind\_ act $\leftarrow$ k 
        \State \Return ind\_ act 
    \end{algorithmic}
\end{algorithm}

\textbf{Evaluation de la complexité en espace}
\begin{itemize}
    \item On ne crée que des variables entières à chaque appel
    \item $ \theta(1) $, et c'est optimal
\end{itemize}


\textbf{Evaluation de la complexité temporelle}
\begin{itemize}
    \item création des variables entières : $ \theta(1) $
    \item Parcours de i à |T|-1 : $ \theta(|T|-1-i) $
    \item Pour n = |T|-1-i la taille de l'entrée : $ \theta(n) $
    \item Etant donné que la comparaison et les affectations sont en temps constant 
    \item La complexité \textbf{ totale } est $\theta(n)$, et c'est optimal (car theta).
\end{itemize}

\subsection*{Question 4}

Invariant de boucle pour cet algorithme : 
\textit{"Tout élément de [i,|T|[ est inférieur ou égal à min\_ act."}

\section*{\textbf{Exercice 2}}

\subsection*{Question 1}

\begin{lstlisting}[label=Python, caption={sontCompatibles}]

def sontCompatibles(i:int, j:int, datedebut:Array[int], datefin:Array[int]) -> bool : 
    int fin_i = datefin[i]
    int debut_j = datedebut[j]
    # Par symetrie des roles de i et j on peut supposer que i est le premier evenement (entre autre i < j)
    if fin_i < debut_j : 
        return true
    else:
        return false

\end{lstlisting}

\subsection*{Question 2}
Supposons que l'on connaisse un évènement sélectionné. 

Dès lors l'ensemble des évènements incompatibles avec ce dernier sont incompatibles pour l'ensemble sélectionné. 

\subsection*{Question 3}

Hypothèse : Choisir en $1^{er}$ élément un évènement de durée minimale.

Ce n'est pas forcément optimal. 
\\
Contre exemple : 

En reprenant l'exemple du cours, si l'on suit cette méthode on peut obtenir \{5, 6, 7, 8\}

Or il existe au moins un sous ensemble de cardinal supérieur : \{ 1, 3, 4, 6, 7\}

\subsection*{Question 4}

Hypothèse : Choisir comme $1^{er}$ élément un évènement de date de fin minimale. 

Cette hypothèse semblerait contribuer à une solution optimale. 
\\
Exemple : 

En reprenant l'exemple du cours, on peut obtenir en appliquant cette hypothèse le sous-ensemble de cardinalité maximale : \{3, 4, 5, 6, 7\}

\subsection*{Question 5}

\textbf{Remarque :}
Pour cette première méthode, on utilise la fonction de l'exercice 1 \textit{IndiceDuMinimum} en ajoutant une comparaison ($ \theta (1) en complexité temporelle $) qui permet de vérifier que le minimum n'a pas déjà été traité avant.
\\ 
Donc on garde une complexité temporelle de $ \theta(n) $ pour cette fonction. 

\begin{algorithm}[H]
    \caption{\textit{Première méthode}}
    \begin{algorithmic}[1]   
        \State E : i, j deux entiers, datedebut et datefin deux tableaux d'entiers
        \State S : un tableau Res de booléen
        \State def Selection\_evt(i:int, j:int, datedebut:int, datefin:int) -> Array[bool] :
        \State \qquad int n $ \leftarrow $ size(datedebut)
        \State \qquad Array[bool] selectionne $ \leftarrow $ New\_Array(n, false)
        \State \qquad Array[bool] visites $ \leftarrow $ New\_Array(n, false)
        \State \qquad int marquage $ \leftarrow $ 0
        \State \qquad While marquage < n:
        \State \qquad \qquad int j $ \leftarrow $ indiceDuMinimum(fin, visites)
        \State \qquad \qquad selectionne[j] $ \leftarrow $ true
        \State \qquad \qquad visites[j] $ \leftarrow $ true
        \State \qquad \qquad marquage $ \leftarrow $ marquage + 1 
        \State \qquad \qquad for k in range(n):
        \State \qquad \qquad \qquad if not visites[k]:
        \State \qquad \qquad \qquad \qquad if not sontCompatibles(k, j, datedebut, datefin):
        \State \qquad \qquad \qquad \qquad \qquad marquage $ \leftarrow $ marquage + 1
        \State \qquad \qquad \qquad \qquad \qquad visites[k] $ \leftarrow $ true
        \State \qquad Libérer la mémoire de visites en fin de programme
        \State \Return selectionne
    \end{algorithmic}
\end{algorithm}

\textbf{Evaluation de la complexité en espace}
\begin{itemize}
    \item 2 tableaux initialisés avec n la taille de l'entrée 
    \item Un des deux tableaux est libéré à la fin de l'exécution 
    \item $ \theta(2n)  \implies \theta(n)$
\end{itemize}

\textbf{Evaluation de le complexité temporelle}
\begin{itemize}
    \item chaque tour coûte $ \theta(n) $ 
    \item Il peut y avoir n tours de boules
    \item \textbf{Pire cas : } $ \theta(n^2) $
\end{itemize}

\subsection*{Question 6}

Il existe une méthode permettant d'améliorer la complexité temporelle de l'algorithme précédent : 
\begin{itemize}
    \item effectuer un tri adapté en $ \theta(n \log (n)) $ selon la fin minimale j. 
    \item Parcourir une unique fois en $ \theta(n) $.
    \item On sélectionne chaque évènement compatible avec le denrier sélectionné.
\end{itemize}
Ainsi on aurait en temps, puis en espace : 
\begin{itemize}
    \item \textbf{Complexité temporelle :} $ \#f(n) = \theta(n \log (n)) + \theta(n) = \theta (n \log (n)) $
    \item \textbf{Complexité en espace : } $ \theta(n) $ (inchangé par rapport à Q5)
\end{itemize}

\textbf{Récriture de l'algorithme : }

% Ecrire l'algorithme amélioré ici : 


\end{document}